\subsection*{One-Stage Sparse BLAS Functionality}


%\linesep

\begin{listing}[H]
\begin{minted}{c++}
using namespace sparseblas;
csr_wrapper<float> A(values, rowptr, colind, shape, nnz);
// auto A = matrix_opt(csr_wrapper<float>(...), allocator);

// scale() overwrites the values of A 
operation_state state(allocator);
scale(policy, state, 2.3, A);
\end{minted}
\caption{Scaling, $A := \alpha A$}
\end{listing}
%\todo[inline]{The table in Section 3 is assigning to a different matrix, but my understanding is we'd want this in-place, like std::linalg?}


%\linesep

\begin{listing}[H]
\begin{minted}{c++}
using namespace sparseblas;
csr_wrapper<float> A(values, rowptr, colind, shape, nnz);

// matrix_opt is opaque and may contain 
// vendor optimization details
auto A_opt = matrix_opt(A, allocator);

// A is const; function returns Inf norm
operation_state state(allocator);
auto inf_nrm = matrix_inf_norm(policy, state, A);
\end{minted}
\caption{Inf Matrix Norm, $\alpha = \|A\|_\text{inf}$.}
\end{listing}


%\linesep

\begin{listing}[H]
\begin{minted}{c++}
using namespace sparseblas;
csr_wrapper<float> A(values, rowptr, colind, shape, nnz);
// auto A = matrix_opt(csr_wrapper<float>(...), allocator);

operation_state state(allocator);
// A is const; function returns Frobenius norm
auto frob_nrm = matrix_frob_norm(policy, state, A);
// or - without passing a policy, resulting in a default:
auto frob_nrm = matrix_frob_norm(state, A);
\end{minted}
\caption{Frobenius Matrix Norm, $\alpha = \|A\|_F$}
\end{listing}



%################################################################


%################################################################

\subsection*{Two-Stage Sparse BLAS Functionality}


\begin{listing}[H]
\begin{minted}{c++}
using namespace spblas;
// csr_wrapper<float> A,B,D filled elsewhwere
// csr_wrapper<float> C(c_nrows, c_ncols);
// auto A_obj = matrix_opt(A, allocator);
// auto B_obj = matrix_opt(B, allocator);
// auto D_obj = matrix_opt(D, allocator);

operation_state state(allocator);
multiply_inspect(policy, state, 
    alpha, A_obj, B_obj, [ beta, D_obj, ] C); // optional
multiply_compute_symbolic(policy, state, 
    alpha, A_obj, B_obj, [ beta, D_obj, ] C);
index_t nnz = state.get_result_nnz();
//T the user allocates the arrays for C
multiply_fill_symbolic(policy, state, 
    alpha, A_obj, B_obj, [ beta, D_obj, ] C);
multiply_numeric(policy, state, 
    alpha, A_obj, B_obj, [ beta, D_obj, ] C);
// csr_wrapper C is now ready to be used
\end{minted}
\caption{Two-stage variant for the sparse matrix multiplication (SpGEMM) 
$C = \alpha \cdot \text{op}(A) \cdot \text{op}(B) [ + \beta \cdot D ]$.}
\end{listing}

\begin{listing}[H]
\begin{minted}{c++}
using namespace spblas;
// csr_wrapper<float> A,B,D filled elsewhwere
// csr_wrapper<float> C(c_nrows, c_ncols);
// auto A_obj = matrix_opt(A, allocator);
// auto B_obj = matrix_opt(B, allocator);
// auto D_obj = matrix_opt(D, allocator);

operation_state state(allocator);
multiply_inspect(policy, state, 
    alpha, A_obj, B_obj, [ beta, D_obj, ] C); // optional
multiply_compute(policy, state, 
    alpha, A_obj, B_obj, [ beta, D_obj, ] C);
index_t nnz = state.get_result_nnz();
//T the user allocates the arrays for C
multiply_fill(policy, state, 
    alpha, A_obj, B_obj, [ beta, D_obj, ] C);
// csr_wrapper C is now ready to be used
\end{minted}
\caption{Three-stage variant for the sparse matrix multiplication (SpGEMM) 
$C = \alpha \cdot \text{op}(A) \cdot \text{op}(B) [ + \beta \cdot D ]$.}
\end{listing}


%\linesep
\begin{listing}[H]
\begin{minted}{c++}
using namespace spblas;
auto A = std::mdspan(raw_x.data(), m, k);
csr_wrapper<float> B(m, n);

operation_state state(allocator);
convert_inspect(policy, state, A, B);
convert_compute(policy, state, A, B);
index_t nnz = state.get_result_nnz();
// allocate the memory for B and put them in B
convert_fill(policy, state, A, B); 
\end{minted}
\caption{Sparse Matrix Format Conversion, $B = \text{sparse}(A)$. Note that this will not remove explicit zeros except for conversion from dense. The above convert\_* functions might accept the predicate function additionally to remove some entries during conversion, which can give some performance benefit without creating a matrix two times.}
\end{listing}


%\linesep

\begin{listing}[H]
\begin{minted}{c++}
using namespace spblas;
csr_wrapper<float> A(values, rowptr, colind, shape, nnz);
// auto A = matrix_opt(csr_wrapper<float>(...), allocator);
// ... likewise for B
csr_wrapper<float> C(m, n);
// auto C = matrix_opt(csr_wrapper<float>(...), allocator);

operation_state state(allocator);
multiply_elementwise_inspect(policy, state, A, B, C); // optional
multiply_elementwise_compute(policy, state, A, B, C);
index_t nnz = state.get_result_nnz();
// allocate C arrays and put in C
multiply_elementwise_fill(policy, state, A, B, C);
// C structure and values are now able to be used
\end{minted}
\caption{Element-wise Multiplication, $C = A~.*B$.}
\end{listing}

%\linesep

\begin{listing}[H]
\begin{minted}{c++}
using namespace spblas;
csr_wrapper<float> A(values, rowptr, colind, shape, nnz);
// auto A = matrix_opt(csr_wrapper<float>(...), allocator);
// ... likewise for B
csr_wrapper<float> C(m, n);
// auto C = matrix_opt(csr_wrapper<float>(...), allocator);

operation_state state(allocator);
add_inspect(policy, state, A, B, C); // optional
add_compute(policy, state, A, B, C);
index_t nnz = state.get_result_nnz();
// allocate C arrays and put in C
add_fill(policy, state, A, B, C);
// C structure and values are now able to be used
\end{minted}
\caption{Sparse Matrix -- Sparse Matrix Addition, $C = A + B$.}
\end{listing}

%\linesep
\begin{listing}[H]
\begin{minted}{c++}
using namespace spblas;
csr_wrapper<float> A_view(values, rowptr, colind, shape, nnz);
csr_wrapper<float> B(values, rowptr, colind, shape, nnz);

auto pred = [](auto i, auto j, auto v) {
              return v > 0;
            };

auto pred = [](auto i, auto j, auto v) {
              return (i < 10 ) && (j < 10);
            };

operation_state state(allocator);
           
// matrix_opt is opaque and may contain vendor optimization details
matrix_opt A(A_view, allocator);
filter_compute(policy, state, A, B, pred);
index_t nnz = state.get_result_nnz();
// the user can allocate the arrays for the output structure
// finally, the output structure can be filled
filter_fill(policy, state, A, B, pred); 
\end{minted}
\caption{Predicate Selection, $B = A.*(A>0)$.}
\end{listing}




%################################################################


%################################################################

\subsection*{Utility Functionality}


